\section{Method}
\label{sec:method}

\noindent\textbf{The measured path.}
Every experiment times the same loop, so that a difference between two runs is
attributable to the factor that was varied and not to a different amount of
work:
%
\[
\begin{aligned}
&\text{receive} \rightarrow \text{replay check} \rightarrow \text{open} \\
&\qquad \rightarrow \text{rekey check} \rightarrow \text{seal}
\rightarrow \text{send}.
\end{aligned}
\]
%
The replay check precedes the cipher deliberately. A replayed packet must not
cost a decryption, which is also why the window's cost lands in $a$.

Packets carry a \protoHeaderBytes-byte header and a \protoTagBytes-byte
authentication tag. The nonce is derived from the epoch and sequence number
through the key schedule rather than carried, as deployed protocols do, and that
is what keeps the header at \protoHeaderBytes\ bytes. The header is
authenticated as associated data, so a tampered sequence number, epoch or
timestamp is rejected along with a tampered payload. With
\protoIPUDPOverhead\ bytes of IP and UDP headers, the largest payload that fits
a \protoMTU-byte MTU is \protoMaxPayload\ bytes, and that is the top of the
sweep.

\noindent\textbf{Statistical discipline.}
The dataset assembled for this report contains \provFragments\ fragments. The
experiments absent from it are \provMissing.

Each point is measured \envReplicates\ times. The replicate count is odd so that
the median is an observed value rather than an average of two. Spread is
reported as a percentile-method bootstrap interval on the median rather than a
standard deviation, because the distribution of a per-packet cost is
right-skewed by scheduling events and a symmetric interval would misrepresent
it \cite{hoefler2015}.

A point whose coefficient of variation exceeds \envCVThreshold\ is marked
unstable. It is recorded rather than dropped: excluded points appear hollow in
the figures, so a reader can see how far they sat from the line.

The sweep order is interleaved and seeded rather than nested. A nested loop
confounds the factor with time. If the machine warms up, or another tenant
arrives halfway through, the whole effect lands on whichever level was running
then. Each replicate of each unit is instead a separate entry in one list,
shuffled with a fixed seed that is recorded with the data.

An anchor configuration is measured before and after every sweep, as the median
of \ctlAnchorReplicates\ units at each end rather than as a single observation.
An anchor built from two single measurements carries the same replicate variance
as any other point and would void sweeps on its own noise. If the two ends
disagree by more than ten per cent the sweep is marked void rather than
reported. In the run described here the anchor moved by \anchorDrift, so the
sweep is \anchorVerdict.

\noindent\textbf{Varying the hardware crypto factor.}
The capability mask that clears AES-NI and the carry-less multiply is read by
libcrypto when it loads, and by the dispatcher of the implementations written
here for the same reason. It cannot be changed inside a running process, so each
unit of the sweep runs in a fresh child process with the mask set or cleared.
The child reports back whether the mask took effect, and a unit whose child
disagrees with the request is discarded. Without that check, a mask that
silently failed would appear in the data as a null result.

Applying the mask to both implementations matters. Masking only OpenSSL would
compare its software path against our hardware path, which measures nothing.

\begin{table}[t]
\caption{The measurement host, as recorded with every dataset fragment.}
\label{tab:environment}
\centering\footnotesize
\input{generated/table_environment}
\end{table}

\noindent\textbf{Host classification.}
A host that cannot hold a measurement is not permitted to write one. Three
classes are enforced in code: \emph{development}, which builds, self-tests and
typesets but never measures; \emph{constrained}, which measures and attaches its
caveats to every dataset; and \emph{measurement}, which adds kernel isolation on
the data-plane cores and at least two homogeneous cores. The host used here is a
\envCores-core \envArch\ machine (\envCPU) running \envOS, with an invariant
counter at \envTSCGHz\,GHz, built by \envCompiler\ against \envOpenSSL. It
classifies as \envHostClass\ in isolation tier \envIsolationTier.
Table~\ref{tab:environment} gives its full description, which travels with every
fragment of the dataset.

\noindent\textbf{Topology.}
Where a second host is available the load generator runs on it and sends across
a physical interface; the data plane binds that interface and forwards
re-encrypted packets back. The generator is driven for a whole sweep over one
control connection rather than over per-unit \texttt{ssh}, which would spend
more time on \texttt{ssh} than on measurement. Without a second host the
generator is a thread on the same machine and the packets never leave it. That
fallback is recorded in every affected point, and its consequences are stated
where they matter, in Section~\ref{sec:latency}.
